% from ega0-16-fr.tex
\oldpage[0\textsubscript{IV-1}]{29}

Pour $n=1$, on obtient le cas particulier~:

\begin{corollary}[16.3.2]
\label{0.16.3.2-fr}
\textup{(``Hauptidealsatz'').}
Pour qu'un idéal premier dans un anneau noethérien soit de hauteur $\leq1$, il
faut et il suffit qu'il soit minimal dans l'ensemble des idéaux premiers
contenant un idéal principal convenable.
\end{corollary}

\begin{corollary}[16.3.3]
\label{0.16.3.3-fr}
\textup{(Artin-Tate).}
Soit $A$ un anneau intègre noethérien. Les conditions suivantes sont
équivalentes~:
\begin{enumerate}
  \item[$a)$] $A$ est un anneau semi-local de dimension $\leq1$.
  \item[$b)$] Il existe $f\neq0$ dans $A$ tel que $A_f$ soit un corps.
\end{enumerate}
\end{corollary}

\begin{proof}
La condition $a)$ équivaut à dire que $A$ n'a qu'un nombre fini d'idéaux
premiers $\mathfrak{m}_i\neq0$, et que ces idéaux sont tous maximaux ($0$ étant
un idéal premier). Comme le produit des $\mathfrak{m}_i$ n'est pas nul, puisque
$A$ est intègre, un élément $f\neq0$ de ce produit appartient à tous les
$\mathfrak{m}_i$, donc $(0)$ est le seul idéal premier de l'anneau intègre
$A_f$, autrement dit $A_f$ est un corps. (On notera que cette partie de la
démonstration n'utilise pas le fait que $A$ est noethérien.) Prouvons
maintenant que $b)$ entraîne $a)$~; l'hypothèse $b)$ signifie que tout idéal
premier $\mathfrak{p}\neq0$ de $A$ contient $f$. Soient $\mathfrak{p}_i$
$(1\leq i\leq n)$ les idéaux premiers minimaux parmi ceux qui contiennent $f$
(qui sont en nombre fini puisque $A/fA$ est noethérien)~; il suffit de prouver
que les $\mathfrak{p}_i$ sont des idéaux \emph{maximaux}, car alors tout idéal
premier $\neq0$ contient nécessairement un des $\mathfrak{p}_i$, donc lui est
égal. Supposons donc qu'il existe un idéal maximal $\mathfrak{m}$ contenant un
des $\mathfrak{p}_i$ et distinct des $\mathfrak{p}_i$~; $\mathfrak{m}$ n'est
donc pas contenu dans la réunion des $\mathfrak{p}_i$ (Bourbaki,
\emph{Alg.~comm.}, chap.~II, \S~1, no~1, prop.~2), autrement dit il existe
$g\in\mathfrak{m}$ n'appartenant à aucun des $\mathfrak{p}_i$. Soit
$\mathfrak{q}$ un des idéaux premiers minimaux parmi ceux qui contiennent
$g$~; il résulte du ``Hauptidealsatz'' (16.3.2) que $\mathfrak{q}$ est de
hauteur 1, donc $\neq0$~; il contient par suite $f$, donc aussi un des
$\mathfrak{p}_i$, et comme $g\in\mathfrak{q}$ et que $g$ n'appartient à aucun
des $\mathfrak{p}_i$, $\mathfrak{q}$ est distinct des $\mathfrak{p}_i$~; il
serait donc de hauteur $\geq2$, ce qui est contradictoire.
\end{proof}

\begin{proposition}[16.3.4]
\label{0.16.3.4-fr}
Soient $A$ un anneau semi-local noethérien, $\mathfrak{r}$ son radical, $M$ un
$A$-module de type fini, $\mathfrak{p}_i$ les éléments de
$\operatorname{Supp}(M)$ tels que
$\dim(A/\mathfrak{p}_i)=\dim(M)$. Alors les $\mathfrak{p}_i$ sont des éléments
minimaux de $\operatorname{Ass}(M)$, et l'on a
$\dim(M/\mathfrak{p}_iM)=\dim(M)$. Pour tout $x\in\mathfrak{r}$, on a
\begin{equation}
\label{0.16.3.4.1-fr}
\tag{16.3.4.1}
  \dim(M/xM)\geq\dim(M)-1.
\end{equation}
En outre, pour que les deux membres de (16.3.4.1) soient égaux, il faut et il
suffit que $x$ n'appartienne à aucun des $\mathfrak{p}_i$.
\end{proposition}

\begin{proof}
La première assertion est immédiate, car les $\mathfrak{p}_i$ sont par
définition des éléments minimaux de $\operatorname{Supp}(M)$ (16.1.6), et ces
derniers sont aussi éléments minimaux de $\operatorname{Ass}(M)$ (Bourbaki,
\emph{Alg.~comm.}, chap.~IV, \S~1, no~4, th.~2). En outre,
$\operatorname{Supp}(M/\mathfrak{p}_iM)$ est l'ensemble des idéaux premiers
contenant $\mathfrak{p}_i$ ($0_{\mathrm I}$, 1.7.5), donc
$\dim(M/\mathfrak{p}_iM)=\dim(A/\mathfrak{p}_i)=\dim(M)$.

Posons $N=M/xM$. Si $x_1,\ldots,x_n$ sont des éléments de $\mathfrak{r}$ tels
que $N/(x_1N+\cdots+x_nN)$ soit de longueur finie, cela signifie que
$M/(xM+x_1M+\cdots+x_nM)$ est de longueur finie, d'où l'inégalité (16.3.4.1)
en vertu de (16.2.3).

Le fait que les deux membres de (16.3.4.1) soient égaux lorsque $x$
n'appartient à aucun des $\mathfrak{p}_i$ résulte de (16.3.4.1) et de
(16.2.3.1). Inversement, si
